\documentclass[10pt,a4paper]{amsart}
\usepackage{amsmath,amssymb,amsthm,graphicx}
% Retained preamble for the Stacks Project chapter snapshot in this directory.
% The chapter's own preamble.tex is a copy of the shared Stacks Project preamble; the block
% below is sliced MECHANICALLY, line for line, from it. Nothing was added, removed or
% redefined; the wrapper class, the xr-hyper/hyperref wiring and the externaldocument list
% are omitted because the assembler supplies the document class.
%
% lines taken from preamble.tex: 3, 5, 9, 11, 12, 156, 157, 158, ... (51 lines total)
%
% package coverage restored from the chapter's own preamble (2026-10-13): xy, xy-2cell

\usepackage[all]{xy}
\xyoption{2cell}
\UseAllTwocells

\newenvironment{reference}{\comment}{\endcomment}
\newenvironment{slogan}{\comment}{\endcomment}
\newenvironment{history}{\comment}{\endcomment}
\theoremstyle{plain}
\newtheorem{theorem}[subsection]{Theorem}
\newtheorem{proposition}[subsection]{Proposition}
\newtheorem{lemma}[subsection]{Lemma}
\theoremstyle{definition}
\newtheorem{definition}[subsection]{Definition}
\newtheorem{example}[subsection]{Example}
\newtheorem{exercise}[subsection]{Exercise}
\newtheorem{situation}[subsection]{Situation}
\theoremstyle{remark}
\newtheorem{remark}[subsection]{Remark}
\newtheorem{remarks}[subsection]{Remarks}
\numberwithin{equation}{subsection}
\def\lim{\mathop{\mathrm{lim}}\nolimits}
\def\colim{\mathop{\mathrm{colim}}\nolimits}
\def\Spec{\mathop{\mathrm{Spec}}}
\def\Hom{\mathop{\mathrm{Hom}}\nolimits}
\def\Ext{\mathop{\mathrm{Ext}}\nolimits}
\def\SheafHom{\mathop{\mathcal{H}\!\mathit{om}}\nolimits}
\def\SheafExt{\mathop{\mathcal{E}\!\mathit{xt}}\nolimits}
\def\Sch{\mathit{Sch}}
\def\Mor{\mathop{\mathrm{Mor}}\nolimits}
\def\Ob{\mathop{\mathrm{Ob}}\nolimits}
\def\Sh{\mathop{\mathit{Sh}}\nolimits}
\def\NL{\mathop{N\!L}\nolimits}
\def\CH{\mathop{\mathrm{CH}}\nolimits}
\def\proetale{{pro\text{-}\acute{e}tale}}
\def\etale{{\acute{e}tale}}
\def\QCoh{\mathit{QCoh}}
\def\Ker{\mathop{\mathrm{Ker}}}
\def\Im{\mathop{\mathrm{Im}}}
\def\Coker{\mathop{\mathrm{Coker}}}
\def\Coim{\mathop{\mathrm{Coim}}}
\DeclareMathSymbol{\boxtimes}{\mathbin}{AMSa}{"02}
\def\QCohstack{\mathcal{QC}\!\mathit{oh}}
\def\Cohstack{\mathcal{C}\!\mathit{oh}}
\def\Spacesstack{\mathcal{S}\!\mathit{paces}}
\def\Quotfunctor{\mathrm{Quot}}
\def\Hilbfunctor{\mathrm{Hilb}}
\def\Curvesstack{\mathcal{C}\!\mathit{urves}}
\def\Polarizedstack{\mathcal{P}\!\mathit{olarized}}
\def\Complexesstack{\mathcal{C}\!\mathit{omplexes}}
\def\Pic{\mathop{\mathrm{Pic}}\nolimits}
\def\Picardstack{\mathcal{P}\!\mathit{ic}}
\def\Picardfunctor{\mathrm{Pic}}
\def\Deformationcategory{\mathcal{D}\!\mathit{ef}}
\usepackage{cleveref}
\crefname{originallemma}{original lemma}{original lemmas}
\Crefname{originallemma}{Original lemma}{Original lemmas}
\makeatletter
\newlabel{categories-definition-cartesian-over-C}{{02XK}{source}{}{sourcepointer-1}{}}
\newlabel{categories-definition-cartesian-over-C@cref}{{[originallemma][0][]02XK}{[1][0][]source}{}{}{}}
\makeatother
\crefname{originallemma}{original lemma}{original lemmas}
\Crefname{originallemma}{Original lemma}{Original lemmas}
\makeatletter
\newlabel{categories-lemma-composition-cartesian}{{02XL}{source}{}{sourcepointer-2}{}}
\newlabel{categories-lemma-composition-cartesian@cref}{{[originallemma][0][]02XL}{[1][0][]source}{}{}{}}
\makeatother
\crefname{originallemma}{original lemma}{original lemmas}
\Crefname{originallemma}{Original lemma}{Original lemmas}
\makeatletter
\newlabel{equation-morphisms-qcoh}{{04U2}{source}{}{sourcepointer-3}{}}
\newlabel{equation-morphisms-qcoh@cref}{{[originallemma][0][]04U2}{[1][0][]source}{}{}{}}
\makeatother
\title{Strongly cartesian morphisms of the quasi-coherent module stack}
\author{The Stacks Project Authors}
\date{stacks-project:examples-stacks.tex:local-snapshot-20261012; GNU FDL 1.2 or later}
\begin{document}
\maketitle
\noindent\textit{Editorial note.} Statement span is lines 332-337 of sources/stacks-examples-stacks/pinned-original-source.tex; proof span is lines 339-372, the complete author proof (32 body lines).\par\medskip
\noindent\textit{Editorial note.} Navigation stubs: 3 cited-but-not-included labels are supplied as non-load-bearing pointers with their pinned tags where the snapshot tag index has them.\par\medskip
\noindent\textit{Editorial note.} extra\_packages is empty on purpose: the retained preamble already loads xy with options and a second declaration clashes in the assembler wrapper.\par\medskip
\noindent\textit{Editorial note.} Text is retained verbatim, including the author's own typos and placeholders.\par\medskip
\noindent\textit{Editorial note.} Queue index 3 of queue\_examples-stacks.json (pb=1548, stmt 332-337, proof 339-372, halfA).\par\medskip
\noindent\textit{Editorial note.} stacks-examples-stacks wave-1 (exs1): fresh adjudication accept (2026-10-14); pool per Chief 09:4x authorization (idle-line throughput over kun's backup note, SMC precedent); examples-stacks.tex 27 labels/20 with proof judged 9 accept / 11 reject; claim\_key v2 stacks:examples-stacks.tex:<label>.\par\medskip
\noindent\textit{Allowed standard prerequisites (editorial).}\begin{itemize}
\item Explicit prerequisite boundary. The cited inputs are prerequisites; the argument in the retained span is not replaced by any cited result.
\end{itemize}\medskip
\ifdefined\hypertarget\hypertarget{sourcepointer-1}{}\fi
\noindent\textit{Source reference.} Original lemma 02XK: sources/stacks-examples-stacks/pinned-original-source.tex, source lines 372--372. Cited by the retained statement or proof but not part of the collected spans; supplied only as a navigation pointer. The argument does not depend on its content.\par\medskip
\ifdefined\hypertarget\hypertarget{sourcepointer-2}{}\fi
\noindent\textit{Source reference.} Original lemma 02XL: sources/stacks-examples-stacks/pinned-original-source.tex, source lines 372--372. Cited by the retained statement or proof but not part of the collected spans; supplied only as a navigation pointer. The argument does not depend on its content.\par\medskip
\ifdefined\hypertarget\hypertarget{sourcepointer-3}{}\fi
\noindent\textit{Source reference.} Original lemma 04U2: sources/stacks-examples-stacks/pinned-original-source.tex, source lines 316--321. Cited by the retained statement or proof but not part of the collected spans; supplied only as a navigation pointer. The argument does not depend on its content.\par\medskip
\begin{lemma}
\label{lemma-quasi-coherent-strongly-cartesian}
A morphism $(f, \varphi) : (Y, \mathcal{G}) \to (X, \mathcal{F})$
of $\QCohstack$ is strongly cartesian if and only if the
map $\varphi$ induces an isomorphism $f^*\mathcal{F} \to \mathcal{G}$.
\end{lemma}

\begin{proof}
Let $(X, \mathcal{F}) \in \Ob(\QCohstack)$.
Let $f : Y \to X$ be a morphism of $(\Sch/S)_{fppf}$.
Note that there is a canonical $f$-map $c : \mathcal{F} \to f^*\mathcal{F}$
and hence we get a morphism
$(f, c) : (Y, f^*\mathcal{F}) \to (X, \mathcal{F})$.
We claim that $(f, c)$ is strongly cartesian.
Namely, for any object $(Z, \mathcal{H})$ of $\QCohstack$ we have
\begin{align*}
\Mor_{\QCohstack}((Z, \mathcal{H}), (Y, f^*\mathcal{F}))
& =
\coprod\nolimits_{g \in \Mor_S(Z, Y)}
\Mor_{\QCoh(\mathcal{O}_Z)}(g^*f^*\mathcal{F}, \mathcal{H}) \\
& =
\coprod\nolimits_{g \in \Mor_S(Z, Y)}
\Mor_{\QCoh(\mathcal{O}_Z)}((f \circ g)^*\mathcal{F}, \mathcal{H}) \\
& =
\Mor_{\QCohstack}((Z, \mathcal{H}), (X, \mathcal{F}))
\times_{\Mor_S(Z, X)} \Mor_S(Z, Y)
\end{align*}
where we have used Equation (\ref{equation-morphisms-qcoh}) twice.
This proves that the condition of
Categories, Definition \ref{categories-definition-cartesian-over-C}
holds for $(f, c)$, and hence our claim is true. Now by
Categories, Lemma \ref{categories-lemma-composition-cartesian}
we see that isomorphisms are strongly cartesian and
compositions of strongly cartesian morphisms are strongly cartesian
which proves the ``if'' part of the lemma. For the converse, note
that given $(X, \mathcal{F})$ and $f : Y \to X$, if there exists a
strongly cartesian morphism lifting $f$ with target $(X, \mathcal{F})$
then it has to be isomorphic to $(f, c)$ (see discussion following
Categories, Definition \ref{categories-definition-cartesian-over-C}).
Hence the "only if" part of the lemma holds.
\end{proof}

\end{document}
